\lecture{02}{Анализ алгоритмов оптимизации в непрерывном времени}

\sourcevideo
    {https://www.youtube.com/playlist?list=PLOzRYVm0a65fhKDS8cYIC7YCuVGJ4FOJx}
    {Week 1: Lecture 2: Analyzing Optimization Algorithms
    in Continuous Time Domain}

Вычислительные методы оптимизации обычно реализуются в дискретном
времени, однако их непрерывные аналоги часто проще анализировать
средствами теории динамических систем. Такой подход позволяет
исследовать устойчивость оптимального равновесия, оценивать скорость
сходимости и конструировать новые дискретные методы.

\section{От градиентного спуска к градиентному потоку}

Для минимизации дифференцируемой функции
\[
    f\colon\RR^d\to\RR
\]
градиентный спуск имеет вид
\[
    x^{k+1}
    =
    x^k-\eta\nabla f(x^k),
\]
где \(\eta>0\) --- шаг метода. Перепишем обновление как
\[
    \frac{x^{k+1}-x^k}{\eta}
    =
    -\nabla f(x^k).
\]

Полагая \(t_k=k\eta\) и переходя к формальному пределу
\(\eta\to0\), получаем дифференциальное уравнение
\[
    \dot{x}(t)
    =
    -\nabla f(x(t)).
\]

\begin{definition}[Градиентный поток]
Градиентным потоком функции \(f\) называется динамическая система
\[
    \dot{x}=-\nabla f(x).
\]
\end{definition}

Дискретный градиентный спуск можно рассматривать как явную схему
Эйлера для градиентного потока:
\[
    x^{k+1}
    =
    x^k+\eta\dot{x}^k
    =
    x^k-\eta\nabla f(x^k).
\]

\begin{remark}
Непрерывная система используется прежде всего для анализа и
проектирования метода. Численная реализация по-прежнему требует
дискретизации.
\end{remark}

\begin{definition}[Точка равновесия]
Точка \(x^\star\) называется точкой равновесия системы
\(\dot{x}=F(x)\), если
\[
    F(x^\star)=0.
\]
Для градиентного потока условие равновесия имеет вид
\[
    \nabla f(x^\star)=0.
\]
\end{definition}

Если \(x^\star\) является внутренней точкой локального минимума
дифференцируемой функции, то она является равновесием градиентного
потока. Обратное утверждение без дополнительных предположений
неверно: критическая точка может быть максимумом или седловой
точкой.

\section{Квадратичная функция}

Рассмотрим
\[
    f(x)=\frac12x^2.
\]
Её градиент равен
\[
    \nabla f(x)=x,
\]
поэтому градиентный поток задаётся уравнением
\[
    \dot{x}=-x.
\]

Единственная точка равновесия
\[
    x^\star=0
\]
одновременно является единственным глобальным минимумом функции.

\begin{lemma}
Решение задачи
\[
    \dot{x}=-x,
    \qquad
    x(0)=x_0,
\]
имеет вид
\[
    x(t)=x_0\mathrm{e}^{-t}.
\]
\end{lemma}

\begin{proof}
Если \(x_0=0\), то \(x(t)\equiv0\). При \(x_0\ne0\) решение сохраняет
знак, и разделение переменных даёт
\[
    \frac{\mathrm{d}x}{x}=-\mathrm{d}t.
\]
После интегрирования и учёта начального условия получаем
\[
    \ln\abs{x(t)}
    =
    -t+\ln\abs{x_0},
\]
откуда
\[
    x(t)=x_0\mathrm{e}^{-t}.
\]
\end{proof}

\begin{definition}[Глобальная экспоненциальная устойчивость]
Равновесие \(x^\star\) называется глобально экспоненциально
устойчивым, если существуют \(C\ge1\) и \(\lambda>0\) такие, что для
любого начального состояния и всех \(t\ge0\)
\[
    \norm{x(t)-x^\star}
    \le
    C\mathrm{e}^{-\lambda t}
    \norm{x(0)-x^\star}.
\]
\end{definition}

Для рассматриваемой системы можно взять \(C=1\) и \(\lambda=1\).
Следовательно, минимум \(x^\star=0\) является экспоненциально
устойчивым равновесием.

Значение целевой функции убывает как
\[
    f(x(t))
    =
    \frac12x_0^2\mathrm{e}^{-2t}.
\]
Таким образом, устойчивость равновесия непосредственно задаёт
скорость сходимости алгоритма к оптимальному решению.

\section{Нормированный градиентный поток}

Обычный градиентный поток замедляется вблизи критической точки,
поскольку
\[
    \norm{\nabla f(x)}
    \longrightarrow0
    \qquad
    \text{при }
    x\to x^\star.
\]
Для изменения скорости движения можно нормировать градиент.

\begin{definition}[Нормированный градиентный поток]
Нормированным градиентным потоком называется система
\[
    \dot{x}
    =
    -
    \frac{\nabla f(x)}
         {\norm{\nabla f(x)}},
    \qquad
    \nabla f(x)\ne0.
\]
В критических точках полагается \(\dot{x}=0\).
\end{definition}

Векторное поле может быть разрывным в критических точках. Ниже
рассматриваются абсолютно непрерывные траектории, удовлетворяющие
уравнению почти всюду. Нормирование сохраняет направление
наискорейшего убывания в евклидовой геометрии и делает скорость равной
единице вне критических точек.

Для функции \(f(x)=x^2/2\) получаем
\[
    \dot{x}
    =
    -\frac{x}{\abs{x}}
    =
    -\operatorname{sign}(x),
    \qquad x\ne0.
\]

\begin{lemma}[Конечно-временная сходимость]
Пусть \(\operatorname{sign}(0)=0\). Решение системы
\[
    \dot{x}=-\operatorname{sign}(x),
    \qquad
    x(0)=x_0,
\]
имеет вид
\[
    x(t)
    =
    \operatorname{sign}(x_0)\bigl(\abs{x_0}-t\bigr)_+,
    \qquad
    (s)_+\coloneqq\max\{s,0\}.
\]
Оно достигает равновесия за время
\[
    T=\abs{x_0}.
\]
\end{lemma}

\begin{proof}
При \(x_0>0\) до достижения нуля выполняется \(\dot{x}=-1\), поэтому
\(x(t)=x_0-t\). При \(x_0<0\) выполняется \(\dot{x}=1\), поэтому
\(x(t)=x_0+t\). После достижения нуля траектория остаётся в точке
равновесия.
\end{proof}

Обычный градиентный поток приближается к нулю асимптотически, тогда
как нормированный поток достигает нуля за конечное время. При этом
далеко от оптимума нормирование может уменьшить естественно большой
градиент, поэтому выбор поля должен учитывать геометрию функции и
область движения.

Дискретизация нормированного потока приводит к обновлению
\[
    x^{k+1}
    =
    x^k
    -
    \eta_k
    \frac{\nabla f(x^k)}
         {\norm{\nabla f(x^k)}}.
\]

\begin{algorithm}[H]
\caption{Нормированный градиентный метод}
\label{alg:normalized-gradient}
\begin{algorithmic}[1]

\Require
дифференцируемая функция \(f\),
начальная точка \(x^0\),
последовательность шагов \(\{\eta_k\}_{k\ge0}\)

\For{\(k=0,1,2,\ldots\)}

    \State
    \(g^k\gets\nabla f(x^k)\)

    \If{\(\norm{g^k}=0\)}

        \State
        \Return \(x^k\)

    \EndIf

    \State
    \(x^{k+1}
    \gets
    x^k-\eta_k g^k/\norm{g^k}\)

\EndFor

\end{algorithmic}
\end{algorithm}

Идея изменения масштаба градиента появляется и в практических
методах обучения нейронных сетей. Однако конкретные алгоритмы
AdaGrad и Adam используют более сложное покоординатное масштабирование
и накопленные статистики градиентов.

\section{Квадратичная и четвертичная функции}

Сравним функции
\[
    f_1(x)=\frac12x^2,
    \qquad
    f_2(x)=\frac14x^4.
\]
Обе имеют единственный минимум в точке \(x^\star=0\), но их градиенты
различаются:
\[
    \nabla f_1(x)=x,
    \qquad
    \nabla f_2(x)=x^3.
\]

При \(0<\abs{x}<1\) выполняется
\[
    \abs{x^3}<\abs{x}.
\]
Следовательно, вблизи минимума четвертичная функция создаёт меньший
градиент и более медленное движение.

Градиентные потоки имеют вид
\[
    \dot{x}=-x
    \qquad\text{и}\qquad
    \dot{x}=-x^3.
\]

\begin{lemma}
Решение системы
\[
    \dot{x}=-x^3,
    \qquad
    x(0)=x_0,
\]
равно
\[
    x(t)
    =
    \frac{x_0}
         {\sqrt{1+2x_0^2t}}.
\]
\end{lemma}

\begin{proof}
Для \(x_0\ne0\)
\[
    \frac{\mathrm{d}x}{x^3}=-\mathrm{d}t.
\]
Отсюда
\[
    \frac{1}{2x(t)^2}
    =
    t+\frac{1}{2x_0^2},
\]
что даёт требуемую формулу. При \(x_0=0\) утверждение очевидно.
\end{proof}

Для квадратичной функции расстояние до минимума убывает
экспоненциально:
\[
    \abs{x(t)}=\abs{x_0}\mathrm{e}^{-t}.
\]
Для четвертичной функции оно убывает лишь полиномиально:
\[
    \abs{x(t)}
    =
    \frac{\abs{x_0}}
         {\sqrt{1+2x_0^2t}}
    =
    O(t^{-1/2}).
\]

Следовательно, одинаковое положение минимума не означает одинаковую
скорость оптимизации. Существенна геометрия функции вблизи
оптимального решения.

\section{Сильная выпуклость}

\begin{definition}[Сильная выпуклость]
Дифференцируемая функция \(f\colon\RR^d\to\RR\) называется
\(\mu\)-сильно выпуклой, если существует \(\mu>0\), такое что
\[
    f(y)
    \ge
    f(x)
    +
    \inner{\nabla f(x)}{y-x}
    +
    \frac{\mu}{2}\norm{y-x}^2
\]
для всех \(x,y\in\RR^d\).
\end{definition}

Для дважды дифференцируемой функции это условие эквивалентно
матричному неравенству
\[
    \nabla^2f(x)\succeq\mu I_d
    \qquad
    \text{для всех }x\in\RR^d.
\]

Функция
\[
    f_1(x)=\frac12x^2
\]
является \(1\)-сильно выпуклой, поскольку
\[
    f_1''(x)=1.
\]

Для функции
\[
    f_2(x)=\frac14x^4
\]
имеем
\[
    f_2''(x)=3x^2.
\]
Так как
\[
    \inf_{x\in\RR}f_2''(x)=0,
\]
она выпукла, но не является сильно выпуклой на всей прямой.

Сильная выпуклость препятствует чрезмерному уплощению функции около
минимума. Поэтому для сильно выпуклых задач возможно получать более
сильные оценки скорости сходимости.

\section{Почему одного градиента может быть недостаточно}

Рассмотрим функцию
\[
    f(x)=\abs{x}.
\]
При \(x\ne0\) модуль её производной равен единице:
\[
    \abs{f'(x)}=1.
\]
Следовательно, по величине градиента невозможно определить расстояние
до минимума: точки \(x=1\) и \(x=100\) дают одинаковую величину
градиента.

Аналогичная трудность сохраняется для гладких приближений функции
\(\abs{x}\): без дополнительной связи между градиентом, значением
функции и расстоянием до множества минимумов невозможно в общем
случае гарантировать ускоренную сходимость только на основании
градиентной информации.

\section{Условие Поляка--Лоясиевича}

Сильная выпуклость является достаточным, но не единственным условием
быстрой сходимости.

\begin{definition}[Неравенство Поляка--Лоясиевича]
Функция \(f\) удовлетворяет неравенству Поляка--Лоясиевича с
константой \(\mu>0\), если
\[
    \frac12\norm{\nabla f(x)}^2
    \ge
    \mu\bigl(f(x)-f^\star\bigr)
\]
для всех \(x\), где
\[
    f^\star=\inf_x f(x).
\]
\end{definition}

Это условие связывает величину градиента с текущей ошибкой по
значению функции. В отличие от сильной выпуклости, оно не требует
выпуклости функции.

Рассмотрим невыпуклый пример
\[
    f(x)=x^2+3\sin^2x.
\]
Эта функция не выпукла, поскольку её вторая производная
\[
    f''(x)=2+6\cos(2x)
\]
может быть отрицательной. Вместе с тем
\[
    f(x)\ge0,
    \qquad
    f(x)=0
    \Longleftrightarrow
    x=0,
\]
поэтому она имеет единственный глобальный минимум в нуле. Однако
невыпуклость и единственность минимума сами по себе не доказывают
PL-неравенство: для этого требуется отдельная глобальная оценка
градиента через \(f(x)-f^\star\).

Неравенство Поляка--Лоясиевича и более широкие классы функций,
допускающие ускоренный анализ, рассматриваются в следующих лекциях.

\section{Проектирование методов через динамические системы}

Непрерывная модель позволяет разделить разработку алгоритма на три
этапа. Сначала выбирается векторное поле
\[
    \dot{x}=F(x),
\]
равновесия которого совпадают с решениями оптимизационной задачи.
Затем средствами теории устойчивости исследуется сходимость
траекторий и её скорость. После этого система дискретизируется для
численной реализации.

\begin{algorithm}[H]
\caption{Проектирование оптимизационного метода через непрерывную модель}
\label{alg:continuous-design}
\begin{algorithmic}[1]

\Require
задача \(\min_x f(x)\)

\State
выбрать поле \(F(x)\), для которого
\[
    F(x)=0
    \quad\Longleftrightarrow\quad
    x\in\operatorname*{arg\,min}_{z}f(z)
\]

\State
исследовать устойчивость равновесий системы
\(\dot{x}=F(x)\)

\State
получить оценку скорости сходимости траекторий

\State
выбрать дискретизацию
\[
    x^{k+1}=x^k+\eta_kF(x^k)
\]

\State
исследовать ограничения на шаг \(\eta_k\)

\Ensure
дискретный оптимизационный алгоритм

\end{algorithmic}
\end{algorithm}


Градиентный спуск является дискретизацией градиентного потока
\(\dot{x}=-\nabla f(x)\). Оптимальные критические точки выступают
равновесиями соответствующей динамической системы, а скорость
оптимизации связана с типом их устойчивости. Для квадратичной функции
градиентный поток сходится экспоненциально, тогда как для четвертичной
функции сходимость имеет полиномиальный характер. Нормирование
градиента позволяет изменить скорость движения и в простейшем
скалярном случае обеспечивает конечно-временную сходимость. Возможность
ускорения определяется геометрическими свойствами функции, в
частности сильной выпуклостью или более общим неравенством
Поляка--Лоясиевича.
